\documentclass[12pt]{article}
\usepackage[utf8]{inputenc}
\usepackage[T1]{fontenc}
\usepackage{hyperref}
\title{Safaliba Language Resources for K-12 Curriculum Development}
\author{Compiled by Claude Sonnet 5.5 (claude-sonnet-5-5) for Andr\'as Kornai}
\date{}
\begin{document}
\maketitle

\section{Introduction}
Safaliba (ISO 639-3 saf; Glottocode safa1243) is an Atlantic-Congo language of northern Ghana with a few thousand speakers. Published estimates differ by source (about 5,000 to 10,000). Very little published material turned up in the searches. This dossier lists only what was actually found and does not claim that further materials do not exist.

\section{Primary Educational Resources}
No Safaliba primer, textbook or school curriculum was found in the search results. One source states that Ghanaian language policy does not recognize Safaliba, that few literacy materials exist, and that public schools in Safaliba areas use Gonja and English reading materials from the Ghana Education Service.
\begin{itemize}
\item Voice of America Learning English report on learning to read Safaliba; it mentions Safaliba literacy texts but does not name specific titles: \url{https://learningenglish.voanews.com/a/learn-to-read-safaliba-helps-learn-english/2737208.html}
\end{itemize}

\section{Linguistic and Reference Materials}
\begin{itemize}
\item Wikipedia article on the Safaliba language (background, orthography history, policy context): \url{https://en.wikipedia.org/wiki/Safaliba_language}
\item Omniglot page on Safaliba and its Latin-based writing system, developed in 1996 by Edmund Kuŋi Yakubu and revised by Paul and Jennifer Schaefer between 1999 and 2003: \url{https://omniglot.com/writing/safaliba.htm}
\item Language and numerals page for Safaliba (Northern Region, Ghana): \url{https://lingweb.eva.mpg.de/channumerals/Safaliba.htm}
\item UNESCO World Atlas of Languages entry for Safaliba in Ghana: \url{https://ar.wal.unesco.org/index.php/en/countries/ghana/languages/safaliba}
\end{itemize}
No dictionary or reference grammar was found in the results.

\section{Religious and Community Resources}
No scripture translation or community organisation was confirmed in the search results. The searches suggested that GILLBT (a Ghanaian Bible translation and literacy organisation) and a Safaliba Literacy Facebook page may hold materials, but no URL for either was seen, so none is cited.

\section{Digital and Audio Resources}
\begin{itemize}
\item Font and script support information for Safaliba (Latin): \url{https://www.fontspace.com/languages/info/saf-latn}
\end{itemize}
No audio or video resources were found.

\section{Curriculum Development Context}
Safaliba has an established Latin orthography, but little printed material and no confirmed place in the school system. A native-language curriculum would probably have to be built from scratch, in collaboration with community members and literacy organisations. Gonja and English are the current school languages in the area.

\end{document}